\section{Realizing the primitive metric addition}
\label{sec:primitive-realization}

The velocity loop gives the desired metric at leading order.  We now remove
its oscillating errors to arbitrarily high finite order, correct the remaining
mean error by small increments, and apply
Proposition~\ref{prop:exact-correction}.  The resulting exact immersions retain
the first and second derivatives needed to control the later boundary curves.

\begin{proposition}[Realization of a velocity loop]
\label{prop:primitive-realization}
Let $M$ be a closed smooth surface and let $F:M\to\R^4$ be a smooth immersion
whose second fundamental form is nonzero at every point.  Let $D$ be a disk
with coordinates $x,y$ on a neighborhood of its closure, and let $a$ be a
smooth nonnegative function, positive on $D$ and zero on its exterior.
Suppose the following data are defined smoothly on a neighborhood of the
$2$-jets of $F$ over $\overline D$ and a collar.  For a map $G$ with jets in
that neighborhood, set
\[
 Y=G_y,\qquad C=G_{yy},\qquad
 P=\operatorname{span}(Y,C)^\perp,\qquad
 v=PG_x,\qquad R_0=(|v|^2+a^2)^{1/2}.
\]
Assume that $Y,C$ are independent and $R_0>0$.  Suppose there is a smooth
$2\pi$-periodic vector $V=V(G;\theta)$, depending only on position and the
$2$-jet of $G$, such that
\[
 V\in P,\qquad |V|=R_0,\qquad \avg_\theta V=v.
\]
Assume also that $V-v$ extends smoothly by zero across the boundary of $D$.
These hypotheses hold for the loop of Lemma~\ref{lem:velocity-loop}.

Put $h_*=\gamma(F)+a^2dx^2$, extending the added tensor by zero.  For all
sufficiently small $z>0$ there is a smooth immersion
$\widetilde f_z:M\to\R^4$ with $\gamma(\widetilde f_z)=h_*$.  To state its
derivatives, evaluate the preceding quantities at $G=F$ and write
$X=F_x-v+V$.  Uniformly on $\overline D$, with periodic arguments evaluated
at $\theta=x/z$,
\begin{equation}
\begin{aligned}
 (\widetilde f_z)_x&=X+o(1),&
 (\widetilde f_z)_y&=Y+o(1),\\
 (\widetilde f_z)_{yy}&=C+o(1),&
 (\widetilde f_z)_{xy}&=X_y+o(1),\\
 z(\widetilde f_z)_{xx}&=V_\theta+o(1).
\end{aligned}
\label{eq:realization-asymptotics}
\end{equation}
Here $X_y$ means differentiation with $\theta$ held fixed.  On the exterior
of $D$, $\widetilde f_z$ converges to $F$ in $C^2$.  There is a constant
$c>0$, independent of sufficiently small $z$, such that
\[
 \min_{p\in M}\max_{|u|=1}
       |B_{\widetilde f_z}(u,u)|\ge c,
\]
where the unit vectors use the fixed background tangent norm.
\end{proposition}

All compact neighborhoods of jets used below are chosen inside the open
sets in this statement.  Thus the stated nondegeneracies have positive
uniform margins.  Derivatives of periodic coefficients in $x,y$ will always
hold $\theta$ fixed.  We write $b_z=o_s(c_z)$ when
$|b_z|_{k,s}/c_z\to0$ for every fixed weighted differentiation order $k$.

\subsection{Removing the oscillating coefficients}

Finite-order asymptotic isometric constructions also occur in Poznyak's
planar immersion problem in $\R^3$ \cite[\S2.3, Theorem~5]{Poznyak1973}.
For the present velocity loop, we give the periodic correctors and their
support properties by an explicit recursion.

For a periodic scalar or vector $b$, define
\[
 \avg b=\frac1{2\pi}\int_0^{2\pi}b(\theta)\,d\theta,
 \qquad \Pi b=b-\avg b.
\]
For $\avg b=0$, let $\mathcal I b$ be its unique periodic primitive with
mean zero.  These operations commute with differentiation in the slow
variables $x,y$.

\begin{lemma}[Finite periodic expansion]
\label{lem:finite-periodic-expansion}
Under the hypotheses of Proposition~\ref{prop:primitive-realization}, for
each integer $L\ge2$ there are periodic differential operators
$U_1,\ldots,U_L$ with values in $\R^4$, and symmetric tensor differential
operators $E_1,\ldots,E_{L-1}$, for which
\begin{equation}
 f_z[G]=G+\sum_{j=1}^L z^jU_j(G;x/z)
 \label{eq:finite-ansatz}
\end{equation}
has the exact expansion
\begin{equation}
 \gamma(f_z[G])
   =a^2dx^2+\mathcal P_z(G)
        +\sum_{r=L}^{2L}z^r\mathcal R_r(G;x/z),
 \qquad
 \mathcal P_z(G)=\gamma(G)+\sum_{r=1}^{L-1}z^rE_r(G).
 \label{eq:finite-metric-expansion}
\end{equation}
Each $U_j$ has mean zero.  The operators $U_j,E_r,\mathcal R_r$ are smooth,
supported on $\overline D$, and polynomial in jets above order two, with
smooth coefficients in the $2$-jet and position, and, where appropriate,
$\theta$.  One may take the differential orders of $U_j$ and $E_r$ to be at
most $2j$ and $2r+2$, respectively.
\end{lemma}

\begin{proof}
Put $w=V-v$ and $X=G_x+w$, and take $U_1=\mathcal Iw$.  Orthogonal
decomposition into $P$ and $P^\perp$ gives
\[
 |X|^2=|G_x-v|^2+R_0^2=|G_x|^2+a^2,
 \qquad X\cdot Y=G_x\cdot G_y.
\]
Thus $X,Y$ induce exactly the leading metric
$\gamma(G)+a^2dx^2$.  Moreover $Y\cdot U_1=C\cdot U_1=0$, so
\begin{equation}
 Y\cdot U_{1,y}=(Y\cdot U_1)_y-C\cdot U_1=0.
 \label{eq:first-yy-cancellation}
\end{equation}
This cancels the $yy$ coefficient at order $z$ before any higher corrector
is introduced.

We construct $U_{r+1}$ from $U_1,\ldots,U_r$.  In the formal first
derivatives of \eqref{eq:finite-ansatz}, write
\[
 f_x=X+\sum_{i\ge1}z^i\mathsf A_i,
 \qquad f_y=Y+\sum_{i\ge1}z^i\mathsf B_i,
 \qquad
 \mathsf A_i=U_{i,x}+U_{i+1,\theta},\quad
 \mathsf B_i=U_{i,y}.
\]
Absent terms are set to zero; in particular $\mathsf A_L=U_{L,x}$.
At step $r$, all $\mathsf A_i$ with $i<r$ and all $\mathsf B_i$ with
$i\le r$ are already known.  The new $U=U_{r+1}$ occurs in the $xx,xy$
coefficients at order $z^r$, and in the $yy$ coefficient at order
$z^{r+1}$.  Removing their fluctuations requires
\begin{equation}
 Y\cdot U_\theta=h,\qquad
 Y\cdot U_y=j,\qquad
 \Pi(X\cdot U_\theta)=e,
 \label{eq:periodic-corrector-system}
\end{equation}
where
\begin{align}
 h&=-\Pi\left(X\cdot\mathsf B_r+Y\cdot U_{r,x}
             +\sum_{i=1}^{r-1}\mathsf A_i\cdot\mathsf B_{r-i}\right),
       \label{eq:periodic-rhs-h}\\
 e&=-\Pi\left(X\cdot U_{r,x}
             +\frac12\sum_{i=1}^{r-1}
                        \mathsf A_i\cdot\mathsf A_{r-i}\right),
       \label{eq:periodic-rhs-e}\\
 j&=-\frac12\Pi\sum_{i=1}^{r}
                         \mathsf B_i\cdot\mathsf B_{r+1-i}.
       \label{eq:periodic-rhs-j}
\end{align}
For example, at $r=1$ these are
\[
 h=-\Pi(X\cdot U_{1,y}+Y\cdot U_{1,x}),\qquad
 e=-\Pi(X\cdot U_{1,x}),\qquad
 j=-\tfrac12\Pi|U_{1,y}|^2.
\]
Thus the first correction removes the $xx,xy$ fluctuations at order $z$
and the $yy$ fluctuation at order $z^2$.  Every right side has mean zero.
This is precisely the compatibility required by the first two equations
of \eqref{eq:periodic-corrector-system} when $\avg U=0$.

Set $H=\mathcal I h$.  There is a unique
$T\in\operatorname{span}(Y,C)$ satisfying
\begin{equation}
 Y\cdot T=H,\qquad C\cdot T=H_y-j.
 \label{eq:corrector-span-part}
\end{equation}
It is obtained by inverting the Gram matrix of $Y,C$, and has mean zero.
Since $Y_y=C$, these equations imply
\[
 Y\cdot T_\theta=h,\qquad Y\cdot T_y=j.
\]
If $Z\in P$, then $Y\cdot Z=C\cdot Z=0$ and hence
$Y\cdot Z_y=0$.  The remaining part of $U$ can therefore be written
$Z=\mathcal I W$, where $W\in P$ has mean zero and must solve
\begin{equation}
 \Pi(V\cdot W)=r_0,
 \qquad r_0=e-\Pi(X\cdot T_\theta).
 \label{eq:covariance-equation}
\end{equation}

Here is an explicit solution which remains regular at $a=0$.  Regard the
following matrices as operators on the Euclidean plane $P$:
\[
 \Sigma=\avg\bigl((V-v)\otimes(V-v)\bigr),
 \qquad \mathsf M=\Id_P-\frac{\Sigma}{R_0^2}.
\]
For a unit vector $q\in P$,
\begin{equation}
 R_0^2\langle\mathsf Mq,q\rangle
   =\avg\bigl(R_0^2-(q\cdot V)^2\bigr)+(q\cdot v)^2.
 \label{eq:covariance-positive}
\end{equation}
Both terms are nonnegative.  Equality would force $V$ to take values only
in $\{R_0q,-R_0q\}$.  Continuity on the period circle would make $V$
constant, contrary to the additional equality $q\cdot v=0$.
Consequently $\mathsf M$ is positive definite.  Its inverse is smooth and
uniformly bounded on the compact family under consideration.  At $a=0$,
the loop is constant, $V=v$, and $\mathsf M=\Id_P$.

Define
\begin{equation}
 \begin{aligned}
 m&=\mathsf M^{-1}\avg(Vr_0),\\
 h_1&=r_0+\frac{(V-v)\cdot m}{R_0^2},\\
 W&=\frac{Vh_1-m}{R_0^2}.
 \end{aligned}
 \label{eq:covariance-solve}
\end{equation}
Then $\avg h_1=0$ and
\[
 \avg(Vh_1)=\avg(Vr_0)+\frac{\Sigma m}{R_0^2}=m.
\]
Hence $\avg W=0$, and $|V|^2=R_0^2$ gives
\[
 \Pi(V\cdot W)=h_1-\frac{(V-v)\cdot m}{R_0^2}=r_0.
\]
Thus $U_{r+1}=T+\mathcal I W$ solves
\eqref{eq:periodic-corrector-system} with mean zero.

This construction is triangular: it changes no previously fixed lower
coefficient.  Induction through $r=L-1$, together with
\eqref{eq:first-yy-cancellation}, makes every metric coefficient of orders
$1,\ldots,L-1$ independent of $\theta$.  Define $E_r$ to be that
coefficient.  The remaining finite coefficients give
$\mathcal R_L,\ldots,\mathcal R_{2L}$ in
\eqref{eq:finite-metric-expansion}.

It remains to check the regularity properties needed for the next stage.
The initial coefficient $U_1$ is a smooth function of the $2$-jet.
Differentiation preserves the class of polynomials in higher jets with
smooth low-jet coefficients, as do multiplication, averaging and
$\mathcal I$.  The Gram and covariance inverses depend only on the
$2$-jet.  Equations \eqref{eq:periodic-rhs-h}--\eqref{eq:covariance-solve}
therefore preserve this class.  If $U_i$ has order at most $2i$ for
$i\le r$, the right side $h$ has order at most $2r+1$; the derivative
$H_y$ raises this by at most one.  All other terms have order at most
$2r+2$.  This proves the asserted order of $U_{r+1}$, and the same
coefficient formulas give order at most $2r+2$ for $E_r$.

Finally, $w$ is smoothly zero on the exterior, so $U_1$ is as well.  If
the previous correctors vanish there, all right sides in
\eqref{eq:periodic-rhs-h}--\eqref{eq:periodic-rhs-j} vanish there and the
linear construction gives $U_{r+1}=0$.  The formulas are smooth on the
collar, including its boundary, so this is smooth extension by zero to
every order.  The same reasoning applies to $E_r$ and $\mathcal R_r$.
\end{proof}

\subsection{Correcting the mean and selecting the scales}

The preceding lemma leaves two errors: a finite mean operator
$\mathcal P_z-\gamma$, and a remainder beginning at $z^L$.  We first record
how the latter behaves when the slow map depends on $z$.
Suppose
\begin{equation}
 G\longrightarrow F\text{ in }C^3,
 \qquad j^2G=O_s(1),\qquad s=z^\sigma,
 \qquad 0<\sigma<1.
 \label{eq:slow-map-bounds}
\end{equation}
A jet of order $k>2$ is $O_s(s^{2-k})$.  Since the expansion has finitely
many coefficients, Lemma~\ref{lem:finite-periodic-expansion} supplies an
integer $p_1\ge0$ such that they and their required periodic derivatives
are all $O_s(s^{-p_1})$.  This exponent is independent of the weighted
differentiation order; the constants may depend on that order.  Indeed
each higher-jet monomial loses the sum of the quantities $k-2$ for its
factors, while weighted differentiation of its coefficients loses no
additional power of $s$.

Substitution of $\theta=x/z$ changes these estimates to $O_z(s^{-p_1})$:
a derivative in $\theta$ costs $z^{-1}$, and a slow derivative costs
$s^{-1}\le z^{-1}$.  In particular,
\begin{equation}
 \gamma(f_z[G])-a^2dx^2-\mathcal P_z(G)
       =O_z(z^Ls^{-p_1}).
 \label{eq:fast-remainder-bound}
\end{equation}

For later use, differentiating the finite ansatz gives
\begin{align*}
 f_x&=X+\sum_{j=1}^L z^jU_{j,x}
           +\sum_{j=2}^L z^{j-1}U_{j,\theta},\\
 f_y&=Y+\sum_{j=1}^L z^jU_{j,y},
 &f_{yy}&=C+\sum_{j=1}^L z^jU_{j,yy},\\
 f_{xy}&=X_y+\sum_{j=1}^L z^jU_{j,xy}
           +\sum_{j=2}^L z^{j-1}U_{j,\theta y},\\
 zf_{xx}&=V_\theta+zG_{xx}
       +\sum_{j=1}^L z^{j+1}U_{j,xx}
       +2\sum_{j=1}^L z^jU_{j,x\theta}
       +\sum_{j=2}^L z^{j-1}U_{j,\theta\theta}.
\end{align*}
All coefficients on these right sides are evaluated at $G$ and
$\theta=x/z$.  Each error contains a positive power of $z$ and at most
two further slow derivatives.  Thus
\begin{equation}
 \sigma(p_1+2)<1
 \label{eq:fast-asymptotic-scale}
\end{equation}
makes these errors uniformly small.  The leading terms converge uniformly
to their values at $F$, including $X_y$, because of the $C^3$ convergence
in \eqref{eq:slow-map-bounds}.  The same estimates show
$\|f_z[G]\|_{C^k}=O(z^{-k})$ for each fixed $k$.

We now construct a slow map satisfying \eqref{eq:slow-map-bounds} for
which the mean error is small enough for exact correction.

\begin{proof}[Proof of Proposition~\ref{prop:primitive-realization}]
Choose the integers $R,N$ in Proposition~\ref{prop:exact-correction} for
the fixed positive metric $h_*$.  Fix an integer
\[
 L>N+R+1
\]
and construct the operators in Lemma~\ref{lem:finite-periodic-expansion}.
Obtain $p_1$ as above and let $p$ be the exponent in
Proposition~\ref{prop:small-increment} for $\mathcal P_z$, with its
perturbation parameter equal to $z$.  These exponents may depend on $L$.
The exponent $p$ is independent of the accuracy parameter in that
proposition.

Define
\[
 d_1=\tfrac14,\qquad d_{l+1}=\tfrac65d_l,
\]
and choose a finite $m\ge2$ with $2d_m>N+R+1$.  Next choose
$\sigma>0$ so small that
\begin{equation}
 \sigma<1,\qquad \sigma(p_1+2)<1,\qquad
 \sigma p<\tfrac12,\qquad \sigma<\tfrac14d_1,
 \qquad L-R-\sigma p_1>N.
 \label{eq:slow-scale-choice}
\end{equation}
Set
\[
 \sigma_l=\frac{l\sigma}{m-1}\quad(0\le l\le m-1),\qquad
 s_l=z^{\sigma_{l-1}},\qquad \delta_l=z^{d_l}
       \quad(1\le l\le m).
\]
Thus $s_1=1$ and $s_m=z^\sigma$.  Fix phase data for $F$ from
Lemma~\ref{lem:phase-data}, with reference tensor $h_0=\gamma(F)$.

Start with
\[
 G_1=(1-\delta_1^2)^{1/2}F,
 \qquad
 H_l=\delta_l^{-2}\bigl(\gamma(F)-\mathcal P_z(G_l)\bigr).
\]
The family $G_1$ is uniformly smooth and $\mathcal P_z-\gamma$ begins
with $z$, whereas $\delta_1^2=z^{1/2}$.  Hence $H_1\to h_0$ in every
smooth norm.  We inductively construct $G_l$ so that
\begin{equation}
 j^2G_l,H_l=O_{s_l}(1),\qquad
 G_l\to F\text{ in }C^3,\qquad H_l\to h_0\text{ in }C^0.
 \label{eq:slow-induction}
\end{equation}

For $l<m$, apply Proposition~\ref{prop:small-increment} at $G_l$ with
amplitude $\delta_l$, scales $s_l$ and $\tau=s_{l+1}$, and input tensor
\[
 H_l-(\delta_{l+1}/\delta_l)^2h_0.
\]
This tensor lies in the allowed neighborhood for small $z$.  If
\[
 \kappa=\min\left\{\frac{\sigma}{m-1},\frac12\right\}>0,
\]
the smallness parameter of that proposition satisfies
\[
 \eta=\frac{s_{l+1}}{s_l}+zs_{l+1}^{-p}=O(z^\kappa).
\]
Only now choose its finite accuracy $q_l$ so large that
\[
 d_l+\kappa q_l>2d_{l+1}.
\]
The other error exponent is positive as well:
\[
 3d_l-\sigma_l-2d_{l+1}
      =\tfrac35d_l-\sigma_l>0.
\]
Thus the increment $U$ supplied by that proposition gives
\[
 \mathcal P_z(G_l+U)
      =\gamma(F)-\delta_{l+1}^2h_0
                       +o_{s_{l+1}}(\delta_{l+1}^2).
\]
Set $G_{l+1}=G_l+U$.  Then $H_{l+1}=h_0+o_{s_{l+1}}(1)$.
Furthermore
\[
 |j^2U|_{k,s_{l+1}}\le C_k z^{d_l-\sigma_l}\longrightarrow0,
 \qquad
 \|U\|_{C^3}\le C z^{d_l-2\sigma_l}\longrightarrow0.
\]
Bounds on the old terms improve when the weighted scale decreases, so
\eqref{eq:slow-induction} is preserved.  All choices are finite.  The
choice of $q_l$ affects only the eventual threshold on $z$; it does not
change $p$ or the already fixed scales.

Take $G=G_m$ and $s=s_m$.  The construction gives
\[
 \mathcal P_z(G)-\gamma(F)=O_s(z^{2d_m}).
\]
With $f_z=f_z[G]$, equation~\eqref{eq:fast-remainder-bound} now yields
\begin{equation}
 \|\gamma(f_z)-h_*\|_{C^R}
   \le C\bigl(z^{2d_m-\sigma R}
                    +z^{L-R-\sigma p_1}\bigr)=O(z^N).
 \label{eq:realization-metric-error}
\end{equation}
The first exponent exceeds $N$ because $2d_m>N+R+1$ and $\sigma<1$;
the second was imposed in \eqref{eq:slow-scale-choice}.  We also have
$\|f_z\|_{C^k}=O(z^{-k})$ through $k=R+2$, and the displayed derivative
formulas prove \eqref{eq:realization-asymptotics} for $f_z$ before exact
correction.  On the exterior $f_z=G$, so $f_z\to F$ in $C^3$ there.
Moreover, \eqref{eq:realization-metric-error} and positivity of $h_*$
guarantee that $f_z$ is immersive for small $z$.

To check the remaining hypothesis of
Proposition~\ref{prop:exact-correction}, project $C=F_{yy}$ to the
orthogonal complement of $\operatorname{span}(X,Y)$.  This projection
never vanishes: if $C=bX+cY$, projection to $P$ gives $0=bV$, hence
$b=0$, contrary to the independence of $C,Y$.  Its length consequently
has a positive minimum on the compact position--phase set.  The
convergence of $(f_z)_x,(f_z)_y,(f_z)_{yy}$ implies a uniform positive
lower bound for $|B_{f_z}(\partial_y,\partial_y)|$ on $\overline D$.
The fixed upper bound on the background length of $\partial_y$ converts
this to the unit-vector bound.  On the exterior the same bound follows
from $C^2$ convergence to $F$ and the nonvanishing of $B_F$.

Proposition~\ref{prop:exact-correction} therefore gives smooth
$\widetilde f_z$ inducing $h_*$ with
$\|\widetilde f_z-f_z\|_{C^2}=o(1)$.  This preserves all the derivative
asymptotics, including the scaled $xx$ derivative, and gives the stated
exterior $C^2$ convergence.  The same argument using the $yy$ component
on the disk, and convergence to $F$ on the exterior, gives the positive
second-form lower bound for $\widetilde f_z$.
\end{proof}

The primitive metric is now realized exactly.  The next section uses the
unscaled $yy$ and $xy$ limits together with the exact Gauss equation to
determine the component of $B_{xx}$ along $B_{yy}$.  The scaled $xx$ limit
controls the complementary normal component.  These are precisely the
second-form data needed to preserve the normal and crossing conditions
for the remaining disks.
